% main_report82.tex
% SAMBP Technical Report TR-82/2028
% Cyber-Resilient GOOSE
\documentclass[11pt, a4paper]{article}

\usepackage[T1]{fontenc}
\usepackage[utf8]{inputenc}
\usepackage[margin=25mm]{geometry}
\usepackage{amsmath, amssymb, amsthm}
\usepackage{booktabs, array, multirow, longtable}
\usepackage{graphicx}
\usepackage[table]{xcolor}
\usepackage[hidelinks]{hyperref}
\usepackage{pgfplots}
\pgfplotsset{compat=1.18}
\usepackage{tikz}
\usetikzlibrary{arrows.meta, positioning, calc, fit, backgrounds,
                shapes.geometric, shapes.misc}
\usepackage{caption}
\usepackage{subcaption}
\usepackage{parskip}
\usepackage{siunitx}
\usepackage{pifont}
\usepackage{cite}

\definecolor{sambpblue}{RGB}{0,70,140}
\definecolor{okgreen}{RGB}{0,130,60}
\definecolor{alertred}{RGB}{180,0,0}
\definecolor{warnoran}{RGB}{200,100,0}

\newcommand{\pu}{\,\text{pu}}
\newcommand{\ms}{\,\text{ms}}
\newcommand{\cmark}{\ding{51}}
\newcommand{\xmark}{\ding{55}}

\newtheorem{finding}{Finding}
\newtheorem{remark}{Remark}

\title{\textbf{\textsf{\color{sambpblue}SAMBP Technical Report TR-82/2028}}\\[6pt]
  \Large Cyber-Resilient GOOSE: Anomaly Detection and IEC~62351-6 Security\\[4pt]
  \normalsize Isolation Forest Detector, HMAC-SHA256 Authentication,
  Four Attack Types, \texttt{local\_io\_mode} Fallback: 10 Tests PASS}
\author{SAMBP Research Group\\[2pt]
  \small Smart Adaptive Multi-Bus Protection Research, IIT Madras}
\date{April 2028\quad\textbar\quad Chapter~7 (Digital Twin / Cyber)\quad\textbar\quad
  Standard: IEC~61850-8-1 + IEC~62351-6:2020}

\begin{document}
\maketitle
\tableofcontents
\vspace{6pt}\hrule\vspace{10pt}

\section{Background and Objectives}
\label{sec:background}

\subsection{Motivation}

GOOSE (Generic Object Oriented Substation Event) messaging is the
real-time communication backbone of IEC~61850 substations, used for
protection trip commands, interlocking, and status exchange.
Because GOOSE operates at Layer-2 Ethernet without TCP/IP,
it is fast ($\leq 4\,\ms$) but also vulnerable to cyber attacks:
replay attacks, timing floods, sequence manipulation, and content injection
can cause spurious trips or denial-of-service~\cite{Hong2014}.

The SAMBP Digital Twin adds a cyber-resilient GOOSE layer based on:
(1)~a machine-learning anomaly detector (Isolation Forest) for traffic
analysis, and (2)~IEC~62351-6:2020 HMAC-SHA256 message authentication.

\subsection{Objectives}

\begin{enumerate}
  \item Implement a 5-feature GOOSE anomaly detector with Isolation Forest.
  \item Detect all four attack types: REPLAY, TIMING\_FLOOD,
    SEQUENCE\_SKIP, CONTENT\_ANOMALY.
  \item Achieve false-positive rate $< 1\%$ on normal traffic.
  \item Implement IEC~62351-6 HMAC-SHA256 authentication with key management.
  \item Validate \texttt{local\_io\_mode} fallback after key failures.
  \item Pass 10 unit tests.
\end{enumerate}

\section{Theory and Methodology}
\label{sec:theory}

\subsection{GOOSE Anomaly Detector}

\subsubsection{Feature Vector}

The anomaly detector uses a 5-dimensional feature vector extracted from
each GOOSE message:

\begin{table}[ht]
\centering
\caption{GOOSE anomaly detector: 5-feature vector.}
\label{tab:features}
\begin{tabular}{cll}
\toprule
Feature & Description & Attack sensitivity \\
\midrule
$f_0$ & $|\Delta\text{stnum}|$ (state number change) & REPLAY, SEQUENCE\_SKIP \\
$f_1$ & $|\Delta\text{sqnum} - 1|$ (sqnum deviation from $+1$) & SEQUENCE\_SKIP \\
$f_2$ & $\log(\text{inter-arrival}/\text{retx\_period} + 1)$ & TIMING\_FLOOD \\
$f_3$ & Shannon entropy of dataset bytes & CONTENT\_ANOMALY \\
$f_4$ & $\log(1 + \text{numeric range of dataset})$ & CONTENT\_ANOMALY \\
\bottomrule
\end{tabular}
\end{table}

\subsubsection{Attack Taxonomy and Rules}

Rule-based detection handles structured attacks before the ML layer:

\begin{description}
  \item[REPLAY:] Message age $> 5\,\text{s}$ OR stnum regression
    ($\text{stnum}_k < \text{stnum}_{k-1}$).
  \item[TIMING\_FLOOD:] Inter-arrival time $< 0.5\,\ms$ (below minimum
    GOOSE retransmission interval).
  \item[SEQUENCE\_SKIP:] Discontinuity in stnum or sqnum beyond expected
    transitions.
  \item[CONTENT\_ANOMALY:] Dataset numeric range $> 3\sigma$ above the
    training bound OR Isolation Forest outlier score.
\end{description}

\subsubsection{Isolation Forest}

The Isolation Forest is trained on $n = 263$ normal GOOSE messages
(contamination parameter $= 0.005$, i.e., 0.5\% expected outlier fraction).
The trained forest assigns an anomaly score $s \in [0, 1]$ to each message;
$s > 0.6$ is flagged as anomalous.

\subsection{IEC 62351-6 HMAC-SHA256 Authentication}

The HMAC-SHA256 message authentication code is computed over:
\begin{equation}
  \mathrm{MAC} = \mathrm{HMAC\text{-}SHA256}\bigl(
    \text{gocb\_ref} \,\Vert\,
    \text{app\_id}_{(4B)} \,\Vert\,
    \text{stnum}_{(4B)} \,\Vert\,
    \text{sqnum}_{(4B)} \,\Vert\,
    t_{\mathrm{msg}(8B)} \,\Vert\,
    \text{dataset\_bytes}
  \bigr)
  \label{eq:hmac}
\end{equation}

Key management implements:
\begin{itemize}
  \item \textbf{Active key:} Current signing/verification key.
  \item \textbf{Pending key:} Next key during grace period (both accepted).
  \item \textbf{Grace period:} 60\,s overlap during key rotation.
  \item \textbf{Fallback:} \texttt{local\_io\_mode} asserted after
    \texttt{fail\_threshold} $= 4$ consecutive wrong-key failures.
    In \texttt{local\_io\_mode}, the IED ignores remote GOOSE commands
    and operates from hardwired I/O only.
\end{itemize}

\section{Simulation Results}
\label{sec:results}

\subsection{Anomaly Detector: Normal Traffic}

False-positive rate on 263 normal GOOSE messages:
1 false positive $\to$ FP rate $= 0.38\%$ ($< 1\%$ threshold). \cmark

\subsection{Attack Detection Results}

\begin{table}[ht]
\centering
\caption{Attack detection results.}
\label{tab:attacks}
\begin{tabular}{lccc}
\toprule
Attack type & Detected & Severity & Detection method \\
\midrule
REPLAY         & \cmark Yes & CRITICAL & Rule: stnum regression \\
TIMING\_FLOOD  & \cmark Yes & CRITICAL & Rule: inter-arrival $< 0.5\,\ms$ \\
SEQUENCE\_SKIP & \cmark Yes & HIGH     & Rule: sqnum gap \\
CONTENT\_ANOMALY & \cmark Yes & HIGH   & Isolation Forest \\
\bottomrule
\end{tabular}
\end{table}

\subsection{HMAC-SHA256 Authentication}

\begin{table}[ht]
\centering
\caption{HMAC authentication scenarios.}
\label{tab:hmac}
\begin{tabular}{lccc}
\toprule
Scenario & HMAC valid & Alarm & Action \\
\midrule
Sign + verify (correct key) & True & OK & None \\
Tampered dataset & False & WARN & Discard \\
Unsigned message & False & WARN & Discard \\
\midrule
4 wrong-key failures & --- & CRITICAL & \texttt{local\_io\_mode=True} \\
\bottomrule
\end{tabular}
\end{table}

\subsection{10-Test Pass Summary}

T01--T04: Rule-based detection (one test per attack type). \cmark\\
T05: Isolation Forest false-positive rate $< 1\%$. \cmark\\
T06: HMAC sign/verify (correct key). \cmark\\
T07: HMAC tampered dataset detection. \cmark\\
T08: Unsigned message detection. \cmark\\
T09: Key rotation grace period. \cmark\\
T10: \texttt{local\_io\_mode} after 4 wrong-key failures. \cmark

\textbf{Total: 10/10 PASS.}

\begin{finding}[Hybrid rule + ML detection is effective]
Rule-based detection handles structured attacks (REPLAY, TIMING\_FLOOD,
SEQUENCE\_SKIP) with zero latency, while the Isolation Forest handles
novel CONTENT\_ANOMALY attacks not covered by explicit rules.
False-positive rate $0.38\%$ (263 normal messages, 1 false positive)
is well within the $< 1\%$ threshold.
\end{finding}

\begin{finding}[HMAC-SHA256 authentication prevents content injection]
All tampered GOOSE messages (modified dataset bytes) are correctly
detected by HMAC verification failure. The \texttt{local\_io\_mode}
fallback after 4 wrong-key failures ensures that a key compromise
does not create a denial-of-service condition.
\end{finding}

\section{Conclusion}
\label{sec:conclusion}

TR-82 delivers the cyber-resilient GOOSE module for the SAMBP Digital Twin.
The hybrid rule + Isolation Forest anomaly detector achieves $0.38\%$ false-positive
rate and detects all four attack types. IEC~62351-6 HMAC-SHA256 authentication
with key rotation and \texttt{local\_io\_mode} fallback is validated in 10 unit tests.

\begin{thebibliography}{99}
\bibitem{Hong2014}
J.~Hong \emph{et al.},
``Cyber attack resilient distance protection and GOOSE communication in
digital substations,''
\emph{IEEE Trans.\ Ind.\ Inform.}, vol.~15, no.~8, pp.~4332--4341, 2019.

\bibitem{IEC62351}
IEC~62351-6:2020, \emph{Power Systems Management and Associated Information
Exchange --- Data and Communications Security --- Part~6: Security for
IEC~61850}, IEC, Geneva, 2020.

\bibitem{Liu2011}
Y.~Liu \emph{et al.},
``False data injection attacks against state estimation in electric power grids,''
\emph{ACM Trans.\ Inf.\ Syst.\ Secur.}, vol.~14, no.~1, pp.~1--33, 2011.

\bibitem{Liu2012}
F.~T.~Liu \emph{et al.},
``Isolation-based anomaly detection,''
\emph{ACM Trans.\ Knowl.\ Discov.\ Data}, vol.~6, no.~1, pp.~1--39, 2012.

\bibitem{Kang2015}
B.~Kang \emph{et al.},
``A cyber-physical testbed for power systems: From theory to implementation,''
\emph{IEEE Syst.\ J.}, vol.~12, no.~3, pp.~2913--2922, 2018.
\end{thebibliography}

\end{document}
